\section{Objectifs Initiaux et Leur Évolution}

\subsection{Vision Initiale du Projet}

Au démarrage du projet, en septembre 2025, l'équipe avait défini un cahier des charges ambitieux. L'AGRI-Scout devait naviguer de façon autonome entre les rangées de cultures en utilisant la caméra pour suivre les lignes au sol, utiliser YOLO pour la détection et l'évaluation de la récolte (identifier l'état des cultures, détecter des végétaux, fruits ou signes de maladies), s'alimenter via un régulateur de tension centralisé (module Traco Power) depuis une unique batterie LiPo, et déployer une sonde de sol après chaque rangée pour cartographier les données NPK.

Les fonctionnalités de détection de feuilles, fruits, légumes et pestes via YOLO ont été retirées du périmètre du projet en raison des contraintes de temps, de la complexité d'annotation des données et des limitations de puissance de calcul. Ces fonctionnalités sont documentées comme pistes d'amélioration future à la section~\ref{sec:ameliorations}.

\subsection{Les Contraintes Rencontrées}

La confrontation avec la réalité matérielle du projet a révélé plusieurs obstacles imprévus qui ont nécessité des adaptations profondes de l'architecture initiale.

\subsubsection{Instabilité Électrique}

Sous forte charge logicielle, la Raspberry Pi 4 consomme jusqu'à 5W en pointe. Le régulateur Traco Power, partagé avec les ESC des moteurs, ne pouvait pas absorber les pics de courant lors des accélérations, provoquant des chutes de tension et des redémarrages intempestifs (brown-outs) rendant toute mission autonome impossible.

\subsubsection{Problèmes de Communication Série}

Lors des premiers tests, le protocole de communication par caractères ASCII simples s'est avéré fragile : les perturbations électromagnétiques des moteurs corrompaient les trames, et l'accumulation de données dans le tampon série provoquait des blocages rendant le robot incontrôlable.

\subsubsection{Limites Computationnelles pour la Vision}

L'exécution de modèles YOLO sur la Raspberry Pi 4 n'atteignait que 2--3 images par seconde, insuffisant pour un contrôle réactif. Le projet a donc abandonné YOLO pour la phase de navigation, mais cette technologie reste pertinente pour de futures améliorations (voir section~\ref{sec:ameliorations}).

\subsection{Les Pivots Stratégiques}

Face à ces contraintes, l'équipe a procédé à une série d'adaptations qui ont transformé le prototype initial en un système fonctionnel et fiable.

\begin{table}[H]
\centering
\caption{Récapitulatif de l'évolution des objectifs du projet AGRI-Scout.}
\label{tab:evolution}
\renewcommand{\arraystretch}{1.4}
\begin{tabularx}{\textwidth}{|>{\bfseries}p{3.5cm}|X|X|}
\hline
\rowcolor{headerblue!20}
\textbf{Domaine} & \textbf{Objectif Initial} & \textbf{Solution Finale Adoptée} \\
\hline
Alimentation & Régulateur Traco Power centralisé & Isolation galvanique : powerbank USB-C pour la RPi, LiPo principale pour la traction \\
\hline
\rowcolor{rowalt}
Vision (détection récolte) & YOLO pour détecter végétaux, fruits et pestes & Retiré du périmètre ; segmentation HSV colorimétrique pour marqueurs de récolte \\
\hline
Navigation & Planification de trajectoire globale (Nav2) & Navigation réactive par contrôleur PD sur erreur de centrage \\
\hline
\rowcolor{rowalt}
Protocole Série & Trames ASCII simples & Paquets binaires 5 octets avec octet de début/fin et checksum XOR \\
\hline
Architecture logicielle & Scripts Python directs & ROS 2 Jazzy avec noeuds asynchrones et topics dédiés \\
\hline
\end{tabularx}
\end{table}
